\documentclass[11pt]{article}
\usepackage[margin=2.2cm]{geometry}
\usepackage[T1]{fontenc}
\usepackage[scaled=0.92]{helvet}
\renewcommand{\familydefault}{\sfdefault}
\usepackage{booktabs,fancyvrb}
\fvset{fontsize=\footnotesize}
\renewcommand{\verbatim}{\Verbatim[fontsize=\scriptsize]}
\renewcommand{\endverbatim}{\endVerbatim}
\title{Formal Languages Project Report\\FIRST/FOLLOW, LL(1), SLR(1) and Syntax-Directed Definitions}
\author{Laura Isabel Isaza Londo\~no \and Santiago Cardenas Angarita \and Santiago Meneses Carvajal}
\date{}
\begin{document}
\maketitle

\section{Book sections followed and data structures}
All algorithms come from Aho et al., \emph{Compilers: Principles, Techniques, and Tools}, 2nd ed.
\begin{center}
\begin{tabular}{p{4.6cm}p{5.6cm}p{5cm}}
\toprule
Part & Book reference & File\\
\midrule
FIRST, FOLLOW & Section 4.4.2 (FIRST of strings too) & \texttt{grammar.py}\\
LL(1) condition and table & Section 4.4.3, Algorithm 4.31 & \texttt{ll1.py}\\
Table-driven predictive parser & Section 4.4.4, Algorithm 4.34 & \texttt{ll1.py}\\
LR(0) items, CLOSURE, GOTO & Section 4.6.2 & \texttt{slr.py}\\
SLR(1) ACTION/GOTO table & Section 4.6.4, Algorithm 4.46 & \texttt{slr.py}\\
LR parsing algorithm & Section 4.6.3, Algorithm 4.44 & \texttt{slr.py}\\
SDD, annotated trees & Sections 5.1.1, 5.1.2 & \texttt{sdd.py}, \texttt{sdd\_example.py}\\
Dependency graph, order & Sections 5.2.1, 5.2.2 & \texttt{sdd.py}\\
S- and L-attributed & Sections 5.2.3, 5.2.4 & Section 3.4 of this report\\
\bottomrule
\end{tabular}
\end{center}

\paragraph{Data structures.}
A grammar is a list of productions \texttt{(head, body)}, with the body a tuple of symbols
(the empty tuple is $\varepsilon$). FIRST and FOLLOW are dictionaries from nonterminal to a Python set.
The LL(1) table is a dictionary \texttt{(A, a) -> production number}; a second write with a different
production is recorded as a conflict. An LR(0) item is a pair \texttt{(production, dot)} and a state is a
frozen set of items. ACTION and GOTO are dictionaries indexed by \texttt{(state, symbol)}. The parse tree is made
of \texttt{Node} objects (symbol, production, children, attributes). An SDD is data: a dictionary from each
production to a list of \texttt{Rule} objects (defined attribute, used attributes, function).

\section{Items 1 and 2: our own grammar}
Grammar (nested lists):
\begin{verbatim}
S -> [L] a
L -> SM
M -> ,SM e
\end{verbatim}
It generates strings such as \texttt{a}, \texttt{[a]}, \texttt{[a,a]} and \texttt{[[a],a]}.
It is both LL(1) and SLR(1), so the program shows the menu.

\subsection{FIRST and FOLLOW (computed by the program)}
\begin{verbatim}
First(L) = {[, a}
First(M) = {,, e}
First(S) = {[, a}
Follow(L) = {]}
Follow(M) = {]}
Follow(S) = {,, ], $}
\end{verbatim}

\subsection{LL(1) table}
\begin{verbatim}
LL(1) predictive parsing table:
   ,         [         ]         a         $         
L            L -> SM             L -> SM             
M  M -> ,SM            M -> e                        
S            S -> [L]            S -> a              
\end{verbatim}

No cell has two productions, so the grammar is LL(1) (Section 4.4.3).

\subsection{LR(0) automaton and SLR(1) table}
In the automaton the production $S'\to S$ is added by the program. Production numbers used in the table:
0: $S\to[L]$, 1: $S\to a$, 2: $L\to SM$, 3: $M\to,SM$, 4: $M\to\varepsilon$.
\begin{verbatim}
LR(0) automaton (canonical collection):
I0:
    S -> .[L]
    S -> .a
    S' -> .S
    --S--> I1
    --[--> I2
    --a--> I3
I1:
    S' -> S.
I2:
    S -> .[L]
    S -> [.L]
    S -> .a
    L -> .SM
    --L--> I4
    --S--> I5
    --[--> I2
    --a--> I3
I3:
    S -> a.
I4:
    S -> [L.]
    --]--> I6
I5:
    L -> S.M
    M -> .,SM
    M -> .
    --M--> I7
    --,--> I8
I6:
    S -> [L].
I7:
    L -> SM.
I8:
    S -> .[L]
    S -> .a
    M -> ,.SM
    --S--> I9
    --[--> I2
    --a--> I3
I9:
    M -> .,SM
    M -> ,S.M
    M -> .
    --M--> I10
    --,--> I8
I10:
    M -> ,SM.
\end{verbatim}
\begin{verbatim}
SLR(1) table:
state  ,      [      ]      a      $      | L    M    S    
0             s2            s3            |           1    
1                                  acc    |                
2             s2            s3            | 4         5    
3      r1            r1            r1     |                
4                    s6                   |                
5      s8            r4                   |      7         
6      r0            r0            r0     |                
7                    r2                   |                
8             s2            s3            |           9    
9      s8            r4                   |      10        
10                   r3                   |                
\end{verbatim}

The table has no conflicts, so the grammar is SLR(1).

\subsection{One string parsed step by step: \texttt{[a,a]}}
LL(1) parser (stack with its top on the left):
\begin{verbatim}
STACK                         INPUT               ACTION
S$                            [a,a]$              S -> [L]
[L]$                          [a,a]$              match [
L]$                           a,a]$               L -> SM
SM]$                          a,a]$               S -> a
aM]$                          a,a]$               match a
M]$                           ,a]$                M -> ,SM
,SM]$                         ,a]$                match ,
SM]$                          a]$                 S -> a
aM]$                          a]$                 match a
M]$                           ]$                  M -> e
]$                            ]$                  match ]
$                             $                   accept
yes
\end{verbatim}

SLR(1) parser:
\begin{verbatim}
STACK                         INPUT               ACTION
0                             [a,a]$              shift 2
0 [ 2                         a,a]$               shift 3
0 [ 2 a 3                     ,a]$                reduce S -> a
0 [ 2 S 5                     ,a]$                shift 8
0 [ 2 S 5 , 8                 a]$                 shift 3
0 [ 2 S 5 , 8 a 3             ]$                  reduce S -> a
0 [ 2 S 5 , 8 S 9             ]$                  reduce M -> e
0 [ 2 S 5 , 8 S 9 M 10        ]$                  reduce M -> ,SM
0 [ 2 S 5 M 7                 ]$                  reduce L -> SM
0 [ 2 L 4                     ]$                  shift 6
0 [ 2 L 4 ] 6                 $                   reduce S -> [L]
0 S 1                         $                   accept
yes
\end{verbatim}

\section{Item 3: SDD}
\subsection{Grammar and attributes}
\begin{verbatim}
S -> L
L -> DN
N -> L e
D -> 0 1 2 3 4 5 6 7 8 9
\end{verbatim}
It generates the non-empty strings of digits. It is LL(1) and SLR(1), so we can parse with either parser
(the evaluator uses the LL(1) parser by default). Our SDD computes the alternating sum of the digits
$d_1-d_2+d_3-\dots$ and the number of digits.
\begin{center}
\begin{tabular}{llllp{6.2cm}}
\toprule
Symbol & Attribute & Kind & Type & Meaning\\
\midrule
D & val & synthesized & int & value of the digit\\
L, N & s & inherited & int ($\pm1$) & sign of the first digit of this subtree\\
L, N & v & synthesized & int & alternating sum of the digits of the subtree\\
L, N & len & synthesized & int & number of digits of the subtree\\
S & v, len & synthesized & int & final results (main attribute: S.v)\\
digit terminals & lexval & synthesized & char & the character itself (given)\\
\bottomrule
\end{tabular}
\end{center}
\subsection{Semantic rules}
\begin{center}
\begin{tabular}{ll}
\toprule
Production & Semantic rules\\
\midrule
$S\to L$ & $L.s=1$;\; $S.v=L.v$;\; $S.len=L.len$\\
$L\to D\,N$ & $N.s=-L.s$;\; $L.v=L.s\cdot D.val+N.v$;\; $L.len=1+N.len$\\
$N\to L$ & $L.s=N.s$;\; $N.v=L.v$;\; $N.len=L.len$\\
$N\to\varepsilon$ & $N.v=0$;\; $N.len=0$\\
$D\to d$ ($d=0,\dots,9$) & $D.val=\texttt{int}(d.lexval)$\\
\bottomrule
\end{tabular}
\end{center}
Each synthesized attribute of the head and each inherited attribute of a body nonterminal is defined exactly once
per production, and $S$ has no inherited attributes (the program checks this with \texttt{check\_sdd}).

\subsection{Step 3: example by hand, string \texttt{1234}}
Annotated parse tree (one node per line; the index after the symbol is the node number; an $\varepsilon$ production has no children):
\begin{verbatim}
S_0   [S -> L]   {v=-2, len=4}
  L_1   [L -> DN]   {s=1, v=-2, len=4}
    D_2   [D -> 1]   {val=1}
      1_3   {lexval='1'}
    N_4   [N -> L]   {s=-1, v=-3, len=3}
      L_5   [L -> DN]   {s=-1, v=-3, len=3}
        D_6   [D -> 2]   {val=2}
          2_7   {lexval='2'}
        N_8   [N -> L]   {s=1, v=-1, len=2}
          L_9   [L -> DN]   {s=1, v=-1, len=2}
            D_10   [D -> 3]   {val=3}
              3_11   {lexval='3'}
            N_12   [N -> L]   {s=-1, v=-4, len=1}
              L_13   [L -> DN]   {s=-1, v=-4, len=1}
                D_14   [D -> 4]   {val=4}
                  4_15   {lexval='4'}
                N_16   [N -> e]   {s=1, v=0, len=0}

\end{verbatim}

Dependency graph: one edge ``used attribute $\to$ defined attribute'' for every rule instance:
\begin{verbatim}
  L_1.v -> S_0.v
  L_1.len -> S_0.len
  L_1.s -> N_4.s
  L_1.s -> L_1.v
  D_2.val -> L_1.v
  N_4.v -> L_1.v
  N_4.len -> L_1.len
  1_3.lexval -> D_2.val
  N_4.s -> L_5.s
  L_5.v -> N_4.v
  L_5.len -> N_4.len
  L_5.s -> N_8.s
  L_5.s -> L_5.v
  D_6.val -> L_5.v
  N_8.v -> L_5.v
  N_8.len -> L_5.len
  2_7.lexval -> D_6.val
  N_8.s -> L_9.s
  L_9.v -> N_8.v
  L_9.len -> N_8.len
  L_9.s -> N_12.s
  L_9.s -> L_9.v
  D_10.val -> L_9.v
  N_12.v -> L_9.v
  N_12.len -> L_9.len
  3_11.lexval -> D_10.val
  N_12.s -> L_13.s
  L_13.v -> N_12.v
  L_13.len -> N_12.len
  L_13.s -> N_16.s
  L_13.s -> L_13.v
  D_14.val -> L_13.v
  N_16.v -> L_13.v
  N_16.len -> L_13.len
  4_15.lexval -> D_14.val

\end{verbatim}

A valid evaluation order (topological order of the graph, produced by the program; the leaf attributes \texttt{lexval} are known from the start and are not listed):
\begin{verbatim}
  1. L_1.s = 1
  2. N_16.v = 0
  3. N_16.len = 0
  4. N_4.s = -1
  5. D_2.val = 1
  6. D_6.val = 2
  7. D_10.val = 3
  8. D_14.val = 4
  9. L_13.len = 1
  10. L_5.s = -1
  11. N_12.len = 1
  12. N_8.s = 1
  13. L_9.len = 2
  14. L_9.s = 1
  15. N_8.len = 2
  16. N_12.s = -1
  17. L_5.len = 3
  18. L_13.s = -1
  19. N_4.len = 3
  20. N_16.s = 1
  21. L_13.v = -4
  22. L_1.len = 4
  23. N_12.v = -4
  24. S_0.len = 4
  25. L_9.v = -1
  26. N_8.v = -1
  27. L_5.v = -3
  28. N_4.v = -3
  29. L_1.v = -2
  30. S_0.v = -2
S.v = -2
\end{verbatim}

Result: $S.v=1-2+3-4=-2$ and $S.len=4$.

\subsection{Step 4: classification}
The SDD is \textbf{L-attributed but not S-attributed}. It is not S-attributed because $L.s$ and $N.s$ are inherited.
It is L-attributed (Section 5.2.4): each inherited attribute depends only on inherited attributes of the head or on attributes
of symbols to its left in the body. $L.s$ in $S\to L$ is a constant; $N.s=-L.s$ in $L\to DN$ uses the head's $s$;
$L.s=N.s$ in $N\to L$ uses the head's $s$. The attributes $D.val$ and $L.v$, $N.v$, $len$ are synthesized, so they have no restriction.
Even so, the evaluator uses the general method (dependency graph and topological order).

\subsection{Step 6: tests}
\begin{center}
\begin{tabular}{lll}
\toprule
String & Expected \texttt{S.v} & Case\\
\midrule
\texttt{7} & 7 & shortest string of the language\\
\texttt{12} & $-1$ & two digits\\
\texttt{1234} & $-2$ & the hand example\\
\texttt{9876543210} & 5 & deep recursion (10 levels of $L$/$N$)\\
\texttt{12a} & \texttt{no} & not in the language\\
\bottomrule
\end{tabular}
\end{center}
Output for \texttt{12a}:
\begin{verbatim}
no
\end{verbatim}

Circularity: in \texttt{sdd\_example.py} we temporarily changed the rule $N.s=-L.s$ of $L\to DN$ into $N.s=N.v$
(\texttt{make\_sdd(circular=True)}). Then $N.s$ depends on $N.v$, which depends on $L.v$ of the child, which depends on $L.s=N.s$.
For a string of one digit there is no cycle, for \texttt{12} there is:
\begin{verbatim}
The SDD is circular: these attributes cannot be ordered:
  S_0.v, L_1.v, N_4.s, N_4.v, L_5.s, L_5.v
\end{verbatim}

\section{Problems found and how we solved them}
\begin{itemize}
\item \textbf{Epsilon productions in FIRST.} FIRST of a string stops at the first symbol that cannot derive $\varepsilon$ and only adds
$\varepsilon$ if all symbols can. We wrote one function for strings and used it in FIRST, FOLLOW and the table.
\item \textbf{Same production twice in one table cell.} When $\varepsilon\in$ FIRST$(\alpha)$ and a terminal is both in FIRST$(\alpha)$ and FOLLOW$(A)$, the same production is written twice in the same cell. That is not a conflict, so we only record a conflict when the production is different.
\item \textbf{End marker.} The program appends \texttt{\$} itself; a string that contains \texttt{\$} is rejected so that the parser cannot accept before reading the whole input.
\item \textbf{Parse tree for the evaluator.} Both parsers build the tree while parsing (LL(1): when a production is expanded; SLR(1): at each reduce). Nodes of $\varepsilon$ productions have no children.
\item \textbf{Augmented grammar.} The production $S'\to S$ is added only inside the SLR(1) code, so the user's grammar is not changed; the accept action is the item $S'\to S\cdot$ on \texttt{\$}.
\item \textbf{Circular SDD.} The topological sort stops with vertices that still have pending inputs; these are reported.
\end{itemize}
\end{document}
